% Generated draft: requires mathematical and source review.
\documentclass[11pt]{article}
\usepackage[margin=25mm]{geometry}
\usepackage{fontspec}
\setmainfont{Latin Modern Roman}
\usepackage{amsmath,amssymb,mathtools,unicode-math}
\setmathfont{Latin Modern Math}
\usepackage{enumitem,needspace,xurl,hyperref}
\setlength{\parindent}{0pt}
\setlength{\parskip}{5pt}
\setlength{\emergencystretch}{3em}
\newcommand{\N}{\mathbb{N}}
\newcommand{\Z}{\mathbb{Z}}
\newcommand{\Q}{\mathbb{Q}}
\newcommand{\R}{\mathbb{R}}
\newcommand{\C}{\mathbb{C}}
\newcommand{\F}{\mathbb{F}}
\newcommand{\A}{\mathbb{A}}
\newcommand{\PP}{\mathbb{P}}
\newcommand{\E}{\mathbb{E}}
\newcommand{\eps}{\varepsilon}
\newcommand{\dd}{\,\mathrm{d}}
\begin{document}
\section*{UnsolvedMath ID: 30003824}
\textbf{Source problem code:} catalogue entry
\subsection{Definitions and mathematical statement}
For a scheme \(X\), write \(D^b(X)=D^b(\operatorname{Coh}(X))\) for the bounded derived category of coherent sheaves, and write \(D^b(k)\) for the bounded derived category of finite-dimensional vector spaces over the ground field \(k\). For \(i\geq 0\), let \(\operatorname{Sym}^i C=C^i/\mathfrak S_i\), with \(\operatorname{Sym}^0 C=\operatorname{Spec}k\).

\textbf{Conditional formulation.} Fix the ground field \(k\) intended by the source. Let \(C\) be a smooth projective curve of genus \(g\), and let \(L\) be a line bundle on \(C\) with \(\deg L=1\). To interpret the displayed indexing, assume \(g\geq 2\); to specify the usual moduli-space interpretation, let \(M=M_C(2,L)\) be the coarse moduli space of stable rank-two vector bundles \(E\) on \(C\) satisfying \(\det E\simeq L\). Here stability means that every line subbundle \(F\subset E\) satisfies \(\deg F<\deg(E)/2=1/2\). The genus restriction and stable coarse-moduli interpretation are supplementary qualifications requiring confirmation, since the supplied statement does not explicitly give them. The ground field is also unresolved; this formulation does not assert the conjecture over every field.

Does \(D^b(M)\) admit a full semiorthogonal decomposition
\[
D^b(M)=\left\langle
\mathcal A_0^{(1)},\mathcal A_0^{(2)},
\mathcal A_1^{(1)},\mathcal A_1^{(2)},\ldots,
\mathcal A_{g-2}^{(1)},\mathcal A_{g-2}^{(2)},
\mathcal B
\right\rangle,
\]
where
\[
\mathcal A_i^{(a)}\simeq D^b(\operatorname{Sym}^i C)
\quad(0\leq i\leq g-2,\ a\in\{1,2\}),
\qquad
\mathcal B\simeq D^b(\operatorname{Sym}^{g-1}C)?
\]
Each component is to be realized as a full admissible triangulated subcategory through a fully faithful exact functor. If the ordered components are \(\mathcal T_1,\ldots,\mathcal T_{2g-1}\), semiorthogonality means
\[
\operatorname{Hom}_{D^b(M)}(V,U[n])=0
\quad\text{for all }j>i,\ U\in\mathcal T_i,\ V\in\mathcal T_j,\ n\in\Z.
\]
Fullness means that these components generate all of \(D^b(M)\) as a triangulated category closed under direct summands, with no additional complementary component. For \(g=2\), the list consists precisely of the two \(i=0\) components followed by \(\mathcal B\).

The two \(i=0\) components are exactly the two displayed copies of \(D^b(k)\). The supplied original wording additionally says ``together with the indicated exceptional'' \(D^b(\C)\) components, which could suggest further components, whereas the cleaned display contains no extra components beyond those already counted. The displayed question above follows the cleaned display; the intended component count requires source review. Cases \(g=0,1\) are not specified by that display.
\subsection{Short English statement}
The question asks whether the derived category of the intended rank-two fixed-determinant bundle moduli space is generated semiorthogonally by two copies of the derived category of each symmetric power \(\operatorname{Sym}^i C\), for \(0\leq i\leq g-2\), and one copy for \(i=g-1\). The ground field, genus range, stability and moduli conventions, and potentially ambiguous exceptional-component count require confirmation. The supplied record labels the problem open; that status is not independently verified here.
\subsection{Sources}
\begin{itemize}
\item Supplied local UnsolvedMath registry record, permanent ID \(30003824\), source problem code \texttt{catalogue entry}, title \emph{Derived Decomposition of Fixed-Determinant Bundle Moduli}, in the \emph{Oberwolfach Reports Open Problems} collection. No individual proposer is supplied; the proposed year is recorded as 2018.
\item Supplied source citation: \emph{Interactions between Algebraic Geometry and Noncommutative Algebra (2019)}, Oberwolfach Reports, DOI 10.4171/owr/2018/24, pp. 1465--1515. Original OWR report: \url{https://doi.org/10.4171/owr/2018/24}. The background describes these reports as documenting mathematical workshops and their open-problem discussions. The citation's 2019 designation is retained alongside the record's 2018 proposed year.
\item Supplied literature source, labeled ``current moduli-space decomposition results'': \url{https://doi.org/10.1017/fms.2023.14}. No authors, article title, or theorem identifiers for this reference are supplied.
\item Literature assessment dated \texttt{2026-08-21}: the supplied curation labels the problem open and reports substantial partial semiorthogonal decompositions containing the expected symmetric-power components, while stating that triviality of the complement needed for the conjectured full decomposition remains unproved. This is supplied dated curation, not an independently verified literature assessment. The background explicitly requests independent verification of time-sensitive status and cited claims.
\item The registry supplies the statement-status label \texttt{corrected\_verified} and says that the cleaned statement was checked against the official Oberwolfach report and preserves its hypotheses, notation, and target. These are supplied verification claims; no source visit or independent verification is asserted in this draft.
\end{itemize}
\end{document}
